\chapter{考察}

\section{結果の解釈}
第3章の主結果は、リスク開示の新規性（\texttt{change\_topk}）が、価格反応（\texttt{car\_mm}）よりも投資家の注目（異常出来高：\texttt{abvol\_mkt}）と強く結びつくことを示した。これは、リスク情報の更新が直ちに株価水準を大きく変化させるというよりも、まず市場参加者の注意配分や情報探索行動を通じて現れる可能性を示唆する。すなわち、有価証券報告書の「事業等のリスク」欄の更新は、投資家にとって新情報の可能性がある一方、短期の価格調整としては限定的であり、注目の喚起として観測されやすいという解釈である。

一方で、不確実性指標（\texttt{vol\_abs}）および価格反応（\texttt{car\_mm}）が同一仕様で明確に有意でなかった点は、いくつかの説明が考えられる。第一に、リスク開示はネガティブ情報の告知であるだけでなく、リスクの所在を整理し、投資家の理解を促すという意味で情報環境の改善として評価され得る。したがって、価格反応の符号や大きさは一意ではなく、平均的には相殺されて小さく見える可能性がある。第二に、有価証券報告書の提出日前後には、同日または近接するIR開示や業績情報が同時に存在し得るため、提出日イベントにおける純粋な効果の抽出が難しい。第三に、価格への反映は即時ではなく遅行的である可能性もある。

\section{先行研究との比較}
先行研究では、リスク開示の量や語彙に基づく指標が、株式リスクや市場反応と関連し得ることが示されてきた\cite{Campbell2014}。また、リスク要因の更新（追加・削除）が不確実性と結びつく可能性も報告されている\cite{Lyle2023}。本研究は、従来の頻度や単純な差分では捉えにくい「意味的更新」に焦点を当て、埋め込みベースの類似度で新規性を定量化した点で、これらの研究の流れと整合的である。

ただし、本研究の結果は、価格反応（CAR）よりも注目度（出来高）で関係が強く観測された。これは、リスク情報の更新がまず注意の喚起として現れるという点で、情報処理・注意制約の議論と整合的である。他方、価格反応が弱いことは、リスク開示が既に市場に織り込まれている、あるいは提出日近傍の情報が多く識別が難しい、といった可能性も含む。したがって、先行研究の示す平均的な価格反応と単純に比較するのではなく、測定対象（リスク欄に限定した意味的更新）、市場（日本株）、イベント定義（提出日）といった制度的差異を踏まえた解釈が必要である。

\section{含意}
本研究の含意は、投資家行動と開示実務の両面にある。第一に、リスク欄の「新規性」は、少なくとも短期的には投資家の注目を喚起しやすい。したがって、投資家の情報処理過程を考える上で、開示の量だけでなく内容の更新度合いを捉えることが有用である。第二に、企業にとっては、形式的な記述の反復（いわゆるボイラープレート）よりも、状況変化に応じた具体的な更新が市場参加者の注意を引く可能性が示唆される。これは、規制対応としての開示に留まらず、投資家コミュニケーションとしての開示の質を高める方向性と整合的である。

一方で、注目の喚起が必ずしも価格の下落や不確実性の上昇につながらない可能性も示された。リスク開示の更新は、リスクの顕在化の兆候である場合もあれば、リスク管理や説明責任の強化（理解の促進）である場合もある。したがって、実務的には、更新内容の性質（リスクの増大を示すのか、管理の改善を示すのか）を区別して分析することが望ましい。

\section{限界と課題}
本研究にはいくつかの限界がある。第一に、推定された関係は、厳密な因果効果ではなく提出日周辺の統計的関連として解釈すべきである。新規性は企業の実態変化と同時に生じ得るため、交絡を完全に排除することは難しい。第3章で示したように、事前窓（プレトレンド）でも新規性とアウトカムが相関する場合があり、これは「新規性が高い文書は、提出以前から注目されやすい」といった背景要因の存在を示唆する。

第二に、提出日をイベント日とする設計は制度的に自然である一方、同日・近接日に他の情報（決算や適時開示等）が混在する可能性がある。第三に、新規性指標は埋め込みモデルと前処理（チャンク分割、トークン化等）に依存し、モデル選択や設定変更により値が変化し得る。第四に、カテゴリ別分析は、カテゴリ定義や分類精度、カテゴリ間の相関の影響を受ける。

以上を踏まえ、今後は（i）より強い外生性を持つショックや制度変更を利用した識別、（ii）提出日以外の情報公開タイミングも含めた設計、（iii）新規性の内容（新しいリスクの追加、既存リスクの強調、表現の具体化等）のタイプ分類、（iv）投資家属性やニュース環境を考慮した反応の異質性分析、などにより解釈と外的妥当性を高める必要がある。


